% ============================================================================
% DANE PROJEKTU — TEN PLIK NALEŻY UZUPEŁNIĆ
% ============================================================================
% Zmieniaj wyłącznie tekst znajdujący się między drugą parą nawiasów { }.
% Nie usuwaj poleceń \newcommand ani ich nawiasów.
%
% Przykład:
%   \newcommand{\projecttitle}{Projekt stanowiska pomiarowego}
%
% Znaki specjalne LaTeX-a, m.in. %, &, _ oraz #, należy poprzedzić znakiem \.
% Przykład: 50\%, R\&D, plik\_wynikowy
% ============================================================================

% ETAP PRACY: draft albo final.
% draft — brakujące dane powodują ostrzeżenia, ale PDF nadal powstaje.
% final — brakujące podstawowe dane zatrzymują kompilację.
\newcommand{\projectstatus}{draft}

% === UCZELNIA I PODSTAWOWE DANE PROJEKTU ===
\newcommand{\university}{Politechnika Opolska}
\newcommand{\faculty}{Wydział Elektrotechniki, Automatyki i Informatyki}
\newcommand{\projecttitle}{Tytuł projektu}
% Podtytuł jest opcjonalny. Pozostaw {}, jeżeli nie jest potrzebny.
\newcommand{\projectsubtitle}{}
\newcommand{\coursetitle}{Nazwa przedmiotu}
\newcommand{\projecttype}{Projekt semestralny}
\newcommand{\programme}{Nazwa kierunku}
\newcommand{\studyform}{studia stacjonarne / niestacjonarne}
\newcommand{\semester}{semestr 0}
\newcommand{\instructor}{dr inż. Imię i nazwisko}

% === AUTOR LUB ZESPÓŁ ===
% Znak & oddziela nazwisko od numeru albumu, a \\ kończy wiersz.
% Jeden autor:
%   Imię i nazwisko & 000000 \\
% Zespół:
%   Anna Nowak & 123456 \\
%   Jan Kowalski & 234567 \\
\newcommand{\projectauthors}{%
  Imię i nazwisko & 000000 \\
}
% Wpisz tych samych autorów w jednym wierszu. To pole jest używane wyłącznie
% we właściwościach PDF, dlatego nie stosuje się tu & ani \\.
\newcommand{\projectpdfauthors}{Imię i nazwisko}

% === MIEJSCE I TERMIN ===
\newcommand{\city}{Opole}
\newcommand{\academicyear}{2026/2027}
\newcommand{\submissiondate}{dd.mm.rrrr}

% === LOGO ===
% Ścieżka jest liczona od katalogu głównego projektu.
% Logo można wyłączyć przez ustawienie: \newcommand{\universitylogo}{}
\newcommand{\universitylogo}{assets/branding/logo-weaii.jpeg}
\newcommand{\universitylogowidth}{0.55\textwidth}

% === OPCJONALNE WYKAZY ===
% Dozwolone wartości to wyłącznie true albo false, bez cudzysłowów.
% Wykaz ma sens dopiero wtedy, gdy raport zawiera odpowiednie podpisane elementy.
\newcommand{\showlistoffigures}{false}
\newcommand{\showlistoftables}{false}
\newcommand{\showlistoflistings}{false}

% === OŚWIADCZENIA KOŃCOWE ===
% Obie części są domyślnie włączone. Wartość false można zastosować wyłącznie,
% gdy prowadzący nie wymaga danego oświadczenia albo nie ma ono zastosowania.
\newcommand{\showaideclaration}{true}
\newcommand{\showcontributions}{true}
